\chapter{Revisão bibliográfica}\label{cap:revisao}

Este capítulo apresenta os conceitos utilizados no desenvolvimento do trabalho e discute pesquisas relacionadas ao uso de modelos de linguagem em sistemas de atendimento. São abordados os agentes baseados em modelos de linguagem, as arquiteturas de agente único e multiagente, a comunicação entre agentes e os critérios empregados na avaliação de sistemas conversacionais.

\section{Fundamentação teórica}\label{sec:fundamentacao}

\subsection{Modelos de linguagem em grande escala}
\label{subsec:llm}

Os modelos de linguagem em grande escala são capazes de interpretar instruções e produzir textos a partir do contexto fornecido em sua entrada. Em sistemas de atendimento, esses modelos podem ser utilizados para identificar a intenção do usuário, consultar informações e elaborar respostas adequadas à solicitação recebida.

Diferentemente de um fluxo baseado somente em regras fixas, um sistema apoiado em \glspl{LLM} pode tratar diferentes formas de expressar uma mesma intenção. Essa característica é útil em conversas nas quais o usuário não segue um roteiro predefinido ou fornece as informações de maneira incompleta.

O comportamento do modelo depende das instruções, das informações disponíveis e do histórico incluído no contexto. Uma resposta também pode variar entre execuções, mesmo quando a entrada é semelhante. Por esse motivo, a utilização de \glspl{LLM} em processos de atendimento exige critérios que permitam verificar a correção das respostas e o cumprimento das regras definidas para o serviço.

As informações processadas pelo modelo são representadas por tokens, que correspondem a unidades de texto utilizadas na entrada e na saída. A quantidade de tokens depende do tamanho dos prompts, do histórico da conversa, da base de conhecimento fornecida e das respostas produzidas. Como arquiteturas multiagentes podem realizar mais de uma chamada ao modelo para tratar uma mesma mensagem, o número de chamadas e o consumo de tokens são medidas relevantes para a comparação proposta.

Yang, Xu e Zhang observam que modelos de linguagem podem melhorar o reconhecimento de intenções em comparação com métodos baseados somente em regras ou classificação tradicional. Os autores também destacam o custo computacional e a latência associados ao seu uso~\cite{yang2025}. Embora a latência não seja uma variável avaliada neste trabalho, o estudo demonstra que a utilização de LLMs acrescenta custos operacionais que devem ser considerados na escolha da arquitetura.

\subsection{Agentes baseados em modelos de linguagem}
\label{subsec:agentes-llm}

Um agente baseado em modelo de linguagem combina o modelo com instruções que definem sua função, informações de contexto e mecanismos para executar ações. Em um chatbot de atendimento, essas ações podem incluir classificar uma solicitação, consultar uma base de informações, formular uma resposta ou encaminhar a conversa para outro setor.

As instruções fornecidas ao agente delimitam sua atuação. Um agente de suporte técnico, por exemplo, pode receber regras sobre os testes permitidos, os dados que devem ser solicitados e as situações que exigem transferência para atendimento humano. O histórico da conversa permite utilizar informações fornecidas em mensagens anteriores e evitar que o usuário tenha que repeti-las.

Agentes baseados em modelos de linguagem também podem utilizar mecanismos de memória e recuperação de informações. \textcite{yan2026} destacam que técnicas de geração aumentada por recuperação podem fornecer ao modelo informações obtidas de fontes externas ou de memórias estruturadas. Nesse tipo de abordagem, os conteúdos relevantes são recuperados dinamicamente e inseridos no contexto utilizado para produzir a resposta.

Neste trabalho não será implementado um mecanismo de geração aumentada por recuperação. A base de conhecimento será formada por documentos previamente definidos, contendo regras, informações e procedimentos fictícios de atendimento. No agente único, todo o conteúdo será disponibilizado em um contexto comum. Na arquitetura multiagente, o mesmo conteúdo será dividido conforme a responsabilidade de cada agente.

Moncada et al. utilizam agentes especializados para realizar a ingestão de dados, a classificação de chamados e a geração de respostas em uma aplicação de suporte ao cliente~\cite{moncada2025}. O trabalho mostra que uma tarefa de atendimento pode ser dividida entre componentes com responsabilidades distintas. Essa divisão, entretanto, exige que as informações necessárias sejam transmitidas corretamente entre os agentes.

\subsection{Arquiteturas de agente único e multiagente}
\label{subsec:arquiteturas}

Na arquitetura de agente único, um mesmo agente concentra as funções necessárias ao atendimento. Ele recebe as mensagens, interpreta a solicitação, mantém o histórico, consulta as informações disponíveis e produz a resposta. As regras dos diferentes setores ficam reunidas no contexto desse agente, o que reduz a necessidade de encaminhamentos internos. Por outro lado, o aumento da quantidade de regras e assuntos pode tornar o contexto mais extenso e dificultar a separação das responsabilidades, exigindo que o agente identifique quais instruções devem ser aplicadas em cada situação.

Na arquitetura multiagente, as responsabilidades são divididas entre agentes com funções específicas. Lee, Tseng e Chen~\cite{lee2025} propõem essa organização para sistemas de suporte à operação de telecomunicações, distinguindo um agente principal (\textit{Chief Agent}) e agentes especializados (\textit{Specific Agents}). O agente principal utiliza LangGraph para organizar o fluxo de execução e adota o paradigma ReWOO, que separa o planejamento das ações, a execução das ferramentas e a elaboração da resposta final. Os agentes especializados representam funções de diferentes sistemas e acessam informações por meio de ferramentas, APIs e extração de dados de páginas web. A proposta ilustra a separação entre coordenação e execução de tarefas específicas. Entretanto, os benefícios de eficiência, escalabilidade e flexibilidade discutidos pelos autores não são acompanhados de uma avaliação quantitativa própria que compare as arquiteturas de agente único e multiagente sob condições equivalentes.

Amri et al.~\cite{amri2026} comparam arquiteturas de agente único e multiagente na decisão de encaminhar conversas para atendimento humano. No agente único, um mesmo agente analisa a solicitação e aciona ferramentas de recuperação de informações ou transferência. Na arquitetura multiagente, o processamento inclui a coleta de contexto, a extração de dados do cliente em conversas em andamento e um agente de roteamento, que direciona a solicitação a agentes de perguntas frequentes, reclamações ou transferência. Os experimentos utilizam os modelos GPT-4.1 e Gemini 2.5 Flash sobre o conjunto QiscusCS, composto por 100 diálogos e 4.265 enunciados. Com GPT-4.1, o Macro F1 aumenta de 0,5356 no agente único para 0,6159 no multiagente; com Gemini 2.5 Flash, passa de 0,3333 para 0,6260. Entretanto, a AUC com GPT-4.1 diminui de 0,7017 para 0,6881, enquanto com Gemini aumenta de 0,6053 para 0,6646. Esses resultados mostram que os ganhos dependem do modelo e da métrica considerada, reforçando a necessidade de comparar as arquiteturas sob condições equivalentes.

\textcite{yan2026} classificam as arquiteturas de comunicação de sistemas multiagentes baseados em modelos de linguagem em estruturas planas, hierárquicas, por equipes, de sociedade e híbridas. Nas estruturas hierárquicas, componentes de nível superior coordenam ou direcionam a execução de agentes responsáveis por tarefas específicas. Essa organização facilita a divisão das responsabilidades, mas acrescenta etapas de comunicação e pode concentrar decisões em determinados componentes.

A arquitetura multiagente proposta neste trabalho aproxima-se de uma organização hierárquica com coordenação centralizada. Um orquestrador será responsável por manter o estado da conversa e controlar a execução. O agente de triagem selecionará o primeiro agente especializado. Enquanto o assunto permanecer na mesma área, as mensagens serão encaminhadas ao agente ativo. Quando esse agente identificar uma mudança de assunto, solicitará uma nova triagem, sem escolher diretamente o próximo especialista.

\subsection{Comunicação, orquestração e retenção de contexto}
\label{subsec:orquestracao}

A orquestração corresponde ao controle do fluxo de execução entre os componentes do sistema. Em uma arquitetura multiagente de atendimento, o orquestrador recebe a mensagem, recupera o estado da conversa e aciona a triagem ou o agente especializado que se encontra ativo. O orquestrador não produz a resposta de atendimento, mas controla quais componentes serão executados e registra as decisões tomadas.

A triagem é responsável por classificar o assunto e indicar o agente adequado. Depois dessa decisão, o agente especializado permanece ativo enquanto as mensagens continuarem relacionadas à sua área. Caso identifique uma mudança de assunto, ele retorna ao orquestrador uma solicitação de nova triagem. Dessa forma, os especialistas não realizam transferências diretas entre si.

Em uma conversa com vários turnos, as informações costumam ser fornecidas gradualmente. O usuário pode informar inicialmente que está sem acesso à Internet e somente depois relatar o estado das luzes do equipamento. Se houver troca de agente entre essas mensagens, o novo agente precisa receber as informações necessárias para continuar o atendimento.

Para representar esse comportamento, o sistema deverá manter o histórico das mensagens, os fatos relevantes extraídos da conversa e a identificação do agente ativo. O histórico preservará as mensagens trocadas; os fatos relevantes poderão registrar dados como CEP, condição de cliente e problema relatado; e a identificação do agente ativo indicará qual componente deve receber a mensagem seguinte.

A comunicação entre os agentes não ocorrerá de maneira livre. As trocas serão intermediadas pelo orquestrador e utilizarão saídas estruturadas, contendo a intenção identificada, o destino indicado, a ação solicitada e os fatos que devem permanecer no contexto. Essa estrutura permitirá registrar o caminho percorrido por cada conversa e verificar em qual etapa ocorreu um possível erro.

\textcite{yan2026} analisam os sistemas multiagentes a partir da comunicação estabelecida entre seus componentes. Segundo essa perspectiva, não basta avaliar separadamente a capacidade de cada agente, pois o resultado também depende das mensagens trocadas, do conteúdo transferido e das decisões tomadas durante a execução.

Falhas em um componente podem afetar as etapas seguintes de um sistema multiagente. \textcite{owotogbe2025} discute esse problema ao investigar a robustez de sistemas formados por agentes baseados em modelos de linguagem. Embora a injeção deliberada de falhas não faça parte do escopo deste trabalho, a possibilidade de propagação de erros justifica o registro do fluxo percorrido por cada atendimento.

\subsection{Avaliação de sistemas conversacionais}
\label{subsec:avaliacao-conversacional}

A avaliação de um sistema conversacional pode considerar tanto a resposta entregue ao usuário quanto o processo utilizado para produzi-la. Uma resposta pode estar correta, mas ter sido obtida após encaminhamentos desnecessários, repetição de perguntas ou uso excessivo de chamadas ao modelo.

A correção indica se a resposta respeita as informações e regras disponíveis. A relevância verifica se o conteúdo produzido trata diretamente da solicitação apresentada. A retenção de conhecimento observa se o sistema preserva informações fornecidas em turnos anteriores. A conclusão da tarefa verifica se as necessidades definidas para o cenário foram atendidas.

Neste trabalho, parte da avaliação será realizada com o DeepEval, uma biblioteca destinada à avaliação de aplicações baseadas em modelos de linguagem~\cite{deepeval}. As mensagens do usuário e do assistente serão organizadas como casos de teste conversacionais, aos quais serão aplicadas as métricas inicialmente previstas de relevância dos turnos, retenção de conhecimento e completude da conversa. Essas métricas utilizam um modelo de linguagem como avaliador e fornecem pontuações e justificativas para analisar, respectivamente, a pertinência das respostas ao contexto, a preservação das informações e o atendimento das intenções apresentadas durante a conversa.

Além da avaliação, o DeepEval será utilizado como apoio à preparação dos cenários sintéticos. As descrições e os desfechos sugeridos pela ferramenta serão revisados pelo pesquisador antes da elaboração dos roteiros definitivos, conforme o procedimento apresentado na seção de cenários conversacionais.

Informações objetivas, como a fila selecionada, o número de chamadas ao modelo, o consumo de tokens e a ocorrência de uma transferência, serão verificadas diretamente nos registros da execução. Essa combinação permite complementar a avaliação por LLM com medidas obtidas sem o julgamento do modelo.

Como a avaliação baseada em modelo pode apresentar variações ou interpretações inadequadas, as duas arquiteturas serão analisadas com o mesmo modelo avaliador, as mesmas configurações e os mesmos critérios. Amostras das avaliações automáticas também serão conferidas manualmente.

\section{Trabalhos relacionados}
\label{sec:trabalhos-relacionados}

Os trabalhos analisados abordam diferentes aspectos do uso de modelos de linguagem e sistemas multiagentes. Alguns apresentam aplicações voltadas ao atendimento, enquanto outros tratam do reconhecimento de intenções, da comunicação entre agentes e da robustez dessas arquiteturas.

Moncada et al.~\cite{moncada2025} apresentam uma arquitetura multiagente para automatizar a classificação e a geração de respostas a chamados de suporte técnico. A solução utiliza o framework CrewAI e o modelo GPT-4.1-mini, distribuindo o processamento entre três agentes: o \textit{Reader Agent}, responsável pela leitura e estruturação dos dados de um arquivo CSV; o \textit{Classifier Agent}, que classifica os chamados como incidente, problema, solicitação ou mudança; e o \textit{Analyzer Agent}, que analisa a prioridade, extrai palavras-chave e produz uma resposta contextualizada.

A avaliação utiliza uma amostra de 20 chamados do conjunto \textit{Customer IT Support Ticket}. O classificador acerta 17 casos, alcançando acurácia de 85\%. Entretanto, o desempenho varia entre as categorias: o F1 é de 1,00 para solicitações e mudanças, 0,77 para incidentes e 0,40 para problemas. Os autores associam os erros principalmente à dificuldade de distinguir incidentes de problemas em descrições ambíguas. Para avaliar as respostas, utilizam representações vetoriais produzidas pelo modelo \textit{all-MiniLM-L6-v2} e calculam a similaridade de cosseno com as respostas de referência, obtendo média de 0,756.

O estudo exemplifica a divisão do processamento de chamados entre agentes especializados e combina a avaliação da classificação com a análise semântica das respostas. Contudo, a amostra reduzida limita a abrangência dos resultados, e a similaridade com uma resposta de referência não demonstra, isoladamente, a resolução correta da solicitação. Além disso, o trabalho não apresenta uma comparação controlada com uma arquitetura de agente único nem avalia a continuidade do atendimento em conversas com múltiplos turnos e mudanças de assunto, aspectos contemplados na proposta deste TCC.

No contexto das telecomunicações, \textcite{lee2025} propõem uma arquitetura multiagente para integrar sistemas de suporte à operação por meio de interações em linguagem natural. A organização segue princípios de projeto orientado a domínio, dividindo as responsabilidades entre um agente principal (\textit{Chief Agent}) e agentes especializados (\textit{Specific Agents}), associados a funções dos sistemas operacionais. O agente principal utiliza LangGraph para coordenar o fluxo e adapta o paradigma ReWOO às etapas de planejamento, execução e consolidação dos resultados. Os agentes especializados acessam dados e executam tarefas por meio de ferramentas, extração de informações de páginas web e APIs. A integração com LangServe permite disponibilizar os agentes como serviços acessíveis por interfaces REST.

Os autores discutem benefícios de modularidade, flexibilidade na escolha dos modelos e possibilidade de execução paralela. Entretanto, não apresentam uma avaliação quantitativa própria que permita medir esses benefícios em comparação com uma arquitetura de agente único. Os ganhos de eficiência no uso de tokens e de acurácia mencionados na descrição do ReWOO são atribuídos ao trabalho de referência desse paradigma, não a um experimento conduzido sobre a arquitetura de telecomunicações proposta. Assim, o estudo contribui como referência de organização e integração entre agentes, mas não permite concluir em quais condições essa organização supera uma solução centralizada. 

O reconhecimento de intenções é investigado por \textcite{yang2025}, que propõem uma abordagem em duas etapas para selecionar agentes em uma plataforma de aplicações de inteligência artificial. Na primeira etapa, o sistema utiliza TF-IDF (\textit{Term Frequency--Inverse Document Frequency}), uma técnica que representa textos por pesos atribuídos às palavras, considerando sua frequência no texto e sua raridade no conjunto de documentos. O sistema calcula a similaridade de cosseno entre essas representações da consulta do usuário e dos nomes e descrições dos agentes disponíveis. Quando a similaridade supera o limiar de 0,1, o agente correspondente é selecionado. Caso contrário, o sistema aciona um LLM para interpretar a solicitação considerando o texto de entrada e os metadados dos agentes.

A avaliação utiliza 30 consultas sintéticas geradas com o modelo Qwen2.5-Turbo, também empregado nos experimentos de reconhecimento de intenções. Os autores comparam três configurações: TF-IDF executado no cliente, TF-IDF executado na nuvem e reconhecimento diretamente por LLM. As duas configurações de TF-IDF apresentam acurácia de 43,5\%, com tempos médios de resposta de 1,28~ms no cliente e 14,56~ms na nuvem. O LLM alcança acurácia de 77,2\%, mas apresenta tempo médio de 2.005,94~ms. Os resultados evidenciam um compromisso entre acurácia e tempo de resposta, motivando o uso do método mais leve antes do acionamento do modelo de linguagem.

Entretanto, as tabelas apresentam os resultados dos métodos separadamente, sem informar a acurácia e o tempo médio do fluxo completo em duas etapas. Também não quantificam o consumo de tokens ou o custo computacional da solução combinada. A avaliação com poucas consultas sintéticas limita a abrangência dos resultados e concentra-se na seleção do agente, sem verificar a continuidade do atendimento, a retenção de informações ou a conclusão das tarefas após o encaminhamento.

\textcite{amri2026} comparam arquiteturas de agente único e multiagente na decisão de encaminhar conversas para atendimento humano. A avaliação utiliza o conjunto QiscusCS, composto por 100 diálogos e 4.265 enunciados, anotados por três especialistas do atendimento da empresa Qiscus. Cada enunciado recebe uma indicação sobre a necessidade de intervenção humana. As duas arquiteturas são avaliadas com GPT-4.1 e Gemini 2.5 Flash, utilizando temperatura de 0,7. As métricas incluem precisão, revocação, F1 por classe, Macro F1 e área sob a curva ROC (AUC).

Com GPT-4.1, a arquitetura multiagente alcança Macro F1 de 0,6159, enquanto o agente único obtém 0,5356. Com Gemini 2.5 Flash, os valores são 0,6260 e 0,3333, respectivamente. Entretanto, a superioridade não ocorre em todas as métricas: com GPT-4.1, a AUC diminui de 0,7017 no agente único para 0,6881 no multiagente; com Gemini, aumenta de 0,6053 para 0,6646. Esses resultados indicam que a divisão de responsabilidades pode melhorar o equilíbrio do desempenho entre as classes, mas seus efeitos dependem do modelo e do critério de avaliação.

Os autores também observam uma tendência de transferência mais cedo nas configurações multiagentes, especialmente com Gemini, o que pode encaminhar para humanos solicitações que o chatbot ainda poderia resolver. A análise evidencia, portanto, um compromisso entre identificar a necessidade de intervenção e preservar a capacidade de atendimento automatizado. Como limitações, a avaliação utiliza dados de uma única empresa, e o conjunto de dados e o código não são disponibilizados publicamente, restringindo a reprodução independente dos experimentos. Além disso, as métricas concentram-se na decisão de transferência, sem avaliar diretamente a qualidade das respostas, a conclusão das tarefas ou o consumo de recursos ao longo do atendimento.

A comunicação em sistemas multiagentes baseados em modelos de linguagem é analisada por \textcite{yan2026}, que organizam a literatura a partir de dois níveis complementares. No nível do sistema, examinam as arquiteturas, os objetivos e os protocolos de comunicação. No nível interno, analisam as estratégias de interação, os paradigmas utilizados para formular e interpretar mensagens, os interlocutores e o conteúdo compartilhado. Essa organização permite examinar não apenas quais agentes compõem um sistema, mas também quando, com quem e de que forma eles trocam informações.

A revisão distingue cinco arquiteturas de comunicação: plana, hierárquica, por equipes, de sociedade e híbrida. Nas arquiteturas planas, os agentes interagem sem uma coordenação central, favorecendo a flexibilidade, mas ampliando o esforço de comunicação conforme cresce o número de participantes. Nas hierárquicas, agentes de níveis superiores coordenam a execução dos demais, facilitando a delegação de tarefas, embora possam introduzir gargalos de processamento e atrasos entre camadas. A organização por equipes reúne especialistas em grupos, enquanto a de sociedade considera interações regidas por normas compartilhadas. As arquiteturas híbridas combinam características dessas organizações.

Entre os principais desafios identificados estão o aumento da demanda computacional, a interpretação inadequada das informações compartilhadas, a interoperabilidade entre protocolos e a segurança das mensagens. Os autores também apontam limitações nas avaliações existentes, que frequentemente privilegiam o desempenho individual dos agentes e deixam de medir propriedades do sistema, como eficiência da coordenação, volume de comunicação e robustez diante de mensagens desatualizadas ou incorretas. A contribuição do artigo consiste na classificação e na análise dos mecanismos descritos na literatura, sem apresentar um experimento próprio que compare agente único e multiagente sob condições controladas.

A robustez dos sistemas multiagentes é discutida por \textcite{owotogbe2025}, que apresenta uma proposta de aplicação da engenharia do caos a sistemas baseados em LLMs. Essa abordagem consiste na introdução deliberada e controlada de falhas para identificar vulnerabilidades e avaliar a capacidade de recuperação do sistema. Entre as perturbações consideradas estão falhas de agentes, atrasos e interrupções na comunicação, contenção de recursos e falhas em cascata, nas quais um problema inicial compromete outros componentes.

A pesquisa adota o paradigma de \textit{design science} e organiza o desenvolvimento em três etapas: levantamento da literatura e de ferramentas para identificar modos de falha; construção de um framework de modelagem e experimentação; e validação por estudos de caso e pesquisa-ação com parceiros industriais. O plano de avaliação combina medidas quantitativas, como tempo de resposta, taxas de detecção de falhas, taxas de erro e utilização de CPU e memória, com análises qualitativas de comportamento, observabilidade e impacto no negócio.

Na publicação, os resultados iniciais restringem-se ao relato de uma revisão multivocal da literatura sobre engenharia do caos, enquanto a análise de repositórios e a validação do framework permanecem em desenvolvimento. Portanto, o artigo oferece uma orientação metodológica para investigar falhas e recuperação, mas ainda não apresenta evidências quantitativas da eficácia da abordagem em sistemas multiagentes baseados em LLMs. Sua contribuição está em destacar que a avaliação dessas arquiteturas deve considerar também as interações entre componentes e a propagação de falhas, além da qualidade das respostas produzidas individualmente.
O \autoref{quadro:trabalhos-relacionados} sintetiza as abordagens, os principais resultados ou contribuições e os limites das avaliações dos trabalhos analisados.

\begin{quadro}[htb]
    \centering
    \caption{Síntese dos trabalhos relacionados}
    \label{quadro:trabalhos-relacionados}
    \footnotesize
    \renewcommand{\arraystretch}{1.15}

    \begin{tabularx}{\textwidth}{
        |>{\raggedright\arraybackslash}p{2.5cm}
        |>{\raggedright\arraybackslash}X
        |>{\raggedright\arraybackslash}X
        |>{\raggedright\arraybackslash}X|
    }
        \hline
        \textbf{Trabalho} &
        \textbf{Abordagem} &
        \textbf{Resultados ou contribuições} &
        \textbf{Limites da avaliação} \\
        \hline

        \textcite{moncada2025} &
        Agentes de leitura, classificação e análise de chamados,
        utilizando CrewAI e GPT-4.1-mini &
        Acurácia de 85\% em 20 chamados; similaridade de cosseno
        média de 0,756 nas respostas &
        Amostra reduzida; sem comparação com agente único
        ou avaliação de conversas com múltiplos turnos \\
        \hline

        \textcite{lee2025} &
        Agente principal e agentes especializados integrados
        a sistemas de telecomunicações &
        Descrição da coordenação com LangGraph e ReWOO
        e da integração por ferramentas e APIs &
        Sem avaliação quantitativa própria dos benefícios
        ou comparação controlada entre arquiteturas \\
        \hline

        \textcite{yang2025} &
        Reconhecimento de intenções por TF-IDF,
        com acionamento de LLM quando a similaridade é baixa &
        Acurácia de 43,5\% com TF-IDF e 77,2\% com LLM;
        menor tempo de resposta com TF-IDF &
        30 consultas sintéticas; sem resultados consolidados
        do fluxo combinado ou avaliação após o encaminhamento \\
        \hline

        \textcite{amri2026} &
        Comparação entre agente único e multiagente
        na transferência para atendimento humano &
        Macro F1 superior no multiagente com os dois modelos;
        AUC inferior com GPT-4.1 e superior com Gemini &
        Avaliação concentrada na transferência;
        dados de uma empresa e sem análise do consumo de recursos \\
        \hline

        \textcite{yan2026} &
        Revisão da literatura centrada na comunicação
        em sistemas multiagentes &
        Classificação das arquiteturas e dos mecanismos
        de comunicação; identificação de desafios de avaliação &
        Sem experimento próprio de comparação entre
        agente único e multiagente \\
        \hline

        \textcite{owotogbe2025} &
        Proposta de framework de engenharia do caos
        para sistemas multiagentes baseados em LLMs &
        Delineamento metodológico e plano de avaliação
        de falhas, recuperação e utilização de recursos &
        Pesquisa em andamento; sem resultados quantitativos
        de validação do framework na publicação \\
        \hline
    \end{tabularx}

    \fonteproprioautor
\end{quadro}

\section{Síntese da revisão e lacuna de pesquisa}
\label{sec:sintese-revisao}

Os trabalhos analisados apresentam diferentes aplicações da divisão de responsabilidades entre agentes, incluindo a classificação de chamados, o reconhecimento de intenções, a integração de sistemas e o encaminhamento para atendimento humano. Os resultados experimentais mostram benefícios em determinadas métricas, mas também diferenças de desempenho conforme o modelo, a tarefa e o critério de avaliação. As análises sobre comunicação e robustez, por sua vez, destacam que a coordenação entre componentes introduz custos e possibilidades de falha que precisam ser considerados na avaliação do sistema.

Conforme sintetizado no \autoref{quadro:trabalhos-relacionados}, os estudos possuem escopos distintos e oferecem evidências complementares. Entre eles, o trabalho de \textcite{amri2026} é o que mais se aproxima da comparação proposta, pois avalia agente único e multiagente com os mesmos modelos sobre diálogos de atendimento. Entretanto, suas métricas concentram-se na decisão de transferência para humanos. Os demais trabalhos apresentam avaliações de chamados ou intenções, propostas arquiteturais, uma revisão dos mecanismos de comunicação e um plano de investigação da robustez.

Nesse conjunto de estudos, não foi identificada uma comparação que avalie conjuntamente a qualidade das respostas, a retenção de informações, a conclusão das tarefas, os encaminhamentos entre setores e os recursos consumidos durante conversas com múltiplos turnos. Essa lacuna dificulta compreender se os possíveis benefícios da especialização compensam o esforço adicional de coordenação em atendimentos que envolvem mudanças de assunto e diferentes áreas de responsabilidade.

Para investigar essa questão, este trabalho propõe implementar uma arquitetura de agente único e uma arquitetura multiagente para o atendimento textual de um provedor de Internet. As duas utilizarão o mesmo modelo de linguagem, parâmetros de geração equivalentes, a mesma base de conhecimento e os mesmos cenários sintéticos. A comparação considerará tanto o desenvolvimento e o resultado das conversas quanto o número de chamadas ao modelo, o consumo de tokens e o custo estimado de execução. O procedimento de implementação, os cenários e os critérios de avaliação são apresentados no capítulo seguinte.